\section{The Eucharistic Prayer: Common Grammar}

\subsection{What the Instruction says the prayer is}

Now begins the centre and high point of the whole celebration, namely the Eucharistic Prayer,
that is, a prayer of thanksgiving and sanctification.\endnote{\latin{IGMR} 78: ``Nunc centrum
et culmen totius celebrationis initium habet, ipsa nempe Prex eucharistica, prex scilicet
gratiarum actionis et sanctificationis.'' The paragraph continues that the priest invites the
people to lift their hearts to the Lord in prayer and thanksgiving and associates them with
himself in the prayer he directs in the name of the whole community to God the Father through
Jesus Christ in the Holy Spirit; that the sense of the prayer is that the whole congregation
should join itself with Christ in confessing the great deeds of God and in offering the
sacrifice; and that it requires all to listen to it with reverence and in silence.} Three
elements of that sentence carry weight. The prayer is the high point --- not the institution
narrative within it, and not the elevation. It is a prayer of two operations, thanksgiving and
sanctification, which is why an account that reduces it to a consecratory formula misdescribes
what the Church says she is doing. And it is one prayer: the Instruction speaks of it as a
single whole from the dialogue to the \latin{Amen}.

The Instruction then distinguishes eight principal elements of which the Eucharistic Prayer
consists, and this enumeration is the analytical spine of everything that follows.\endnote{\latin{IGMR} 79 a--h, read in the Latin of the registered reproduction of the 2002
edition and cross-checked in the United States English state of Chapter II. The table
paraphrases the eight lettered paragraphs closely; the Latin of (f) is worth having exactly,
because it is the Instruction's most consequential single sentence about the assembly's part:
``Intendit vero Ecclesia ut fideles non solummodo immaculatam hostiam offerant sed etiam
seipsos offerre discant, et de die in diem consummentur, Christo mediatore, in unitatem cum Deo
et inter se, ut sit tandem Deus omnia in omnibus.''}

{\footnotesize
\begin{longtable}{>{\bfseries\raggedright\arraybackslash}p{0.16\linewidth}>{\raggedright\arraybackslash}p{0.80\linewidth}}
\toprule
\textbf{Element} & \textbf{What \latin{IGMR} 79 says it does}\\
\midrule
\endhead
a. Thanksgiving & Expressed especially in the Preface: the priest, in the name of the whole holy people, glorifies God the Father and gives him thanks for the whole work of salvation or for some particular aspect of it, according to the day, feast, or season.\\
b. Acclamation & The whole congregation, joining itself to the heavenly powers, sings the \latin{Sanctus}; this acclamation constitutes part of the Eucharistic Prayer itself and is uttered by all the people with the priest.\\
c. Epiclesis & By particular invocations the Church implores the power of the Holy Spirit, that the gifts offered by human hands may be consecrated, that is, become the Body and Blood of Christ, and that the spotless victim to be received in Communion may be for the salvation of those who will partake of it.\\
d. Institution narrative and consecration & By the words and actions of Christ that sacrifice is carried out which Christ himself instituted at the Last Supper, when he offered his Body and Blood under the species of bread and wine, gave them to his Apostles to eat and drink, and left them the command to perpetuate the same mystery.\\
e. Anamnesis & Fulfilling the command received from Christ the Lord through the Apostles, the Church keeps the memorial of Christ himself, recalling especially his blessed passion, glorious resurrection, and ascension into heaven.\\
f. Oblation & In this very memorial the Church, and in particular the Church here and now gathered, offers the spotless victim to the Father in the Holy Spirit. The Church intends that the faithful should offer not only the spotless victim but also learn to offer their very selves, and day by day be consummated, through Christ the Mediator, into unity with God and with one another, so that at last God may be all in all.\\
g. Intercessions & These express that the Eucharist is celebrated in communion with the whole Church, both heavenly and earthly, and that the offering is made for her and for all her members, living and dead, who have been called to share in the redemption and salvation acquired by the Body and Blood of Christ.\\
h. Final doxology & The glorification of God is expressed and is confirmed and concluded by the people's acclamation \latin{Amen}.\\
\bottomrule
\end{longtable}}

Element (f) is where the whole doctrine of the assembly's participation is concentrated, and it
should be read with care. The faithful truly offer the spotless victim --- \latin{immaculatam
hostiam offerant} --- and they are to learn to offer themselves as well. Neither clause makes
the assembly a collective celebrant, and the Instruction elsewhere reserves the Eucharistic
Prayer to the priest and requires the rest to listen in silence. The claim is that the offering
of the Church is one offering, made by the priest acting in the person of Christ the Head and
by the baptized in their own manner; the reformed rite states this more explicitly than its
predecessor did, and it is a statement of received doctrine rather than a novelty.

\subsection{The dialogue, the preface, and the \latin{Sanctus} (nn.\ 31--32)}

The prayer opens with the oldest continuously attested dialogue in the Latin liturgy:
\latin{Dominus vobiscum} --- \latin{Et cum spiritu tuo}; \latin{Sursum corda} --- \latin{Habemus
ad Dominum}; \latin{Gratias agamus Domino Deo nostro} --- \latin{Dignum et iustum est}. The
priest then continues the preface with hands extended, and at its end joins his hands and,
together with the people, concludes it by singing or saying the \latin{Sanctus}.\endnote{\latin{Ordo
Missae} 31. The rubric that the priest concludes the preface \latin{una cum populo} is the
Missal's own way of making the point the Instruction states at n.\ 79 b: the \latin{Sanctus} is
not an interruption of the prayer but part of it.} The 2011 English restored ``It is right and
just'' for \latin{Dignum et iustum est} and ``Lord God of hosts'' for \latin{Dominus Deus
Sabaoth}; both are instances of the same principle of closer rendering, and neither is disputed.

Number 32 carries two rubrics that govern the whole section. In all Masses the celebrating
priest is permitted to sing the principal parts of the Eucharistic Prayer, with the music
supplied later in the book; and in the first Eucharistic Prayer, the Roman Canon, the passages
enclosed in parentheses may be omitted.\endnote{\latin{Ordo Missae} 32: ``In Prece eucharistica
prima, seu Canone romano, ea quae inter parentheses includuntur omitti possunt.'' The
parenthesised passages are the second halves of the two saints' lists and the \latin{Per
Christum Dominum nostrum. Amen.} conclusions of several sections. This permission is
postconciliar; in the older Order these passages were obligatory.} That second rubric is the
single largest change to the Roman Canon's manner of use and is treated in the next section.

The prefaces themselves --- more than eighty in the third typical edition --- lie between nn.\ 33
and 82 of the Order of Mass. Their multiplication is one of the reform's stated aims: Paul VI's
constitution says the Eucharistic Prayer was enriched with an abundance of prefaces, taken
either from the more ancient tradition of the Roman Church or newly composed, so that
particular aspects of the mystery of salvation might be more clearly disclosed and richer
motives of thanksgiving supplied.\endnote{Paul VI, \latin{Missale Romanum}: ``\ldots Precatio
eucharistica aucta est copia praefationum, vel ex antiquiore Romanae Ecclesiae traditione
sumptarum, vel nunc primum compositarum \ldots'' The two categories --- recovered and newly
composed --- are the constitution's own and should not be conflated in either direction.}

\subsection{The words of consecration and the removal of \latin{mysterium fidei}}

Paul VI's constitution records a decision that shaped every one of the four prayers, and it is
worth quoting because it is the most exact statement anywhere of what was done and why:\endnote{Paul VI, apostolic constitution \latin{Missale Romanum} (3 April 1969), Latin as
reprinted at the head of the third typical edition and read publication-locally 2026-07-25 in
the registered reproduction. The passage states three distinct decisions: that three new Canons
be added; that the dominical words be one and the same in every Canon formula, for pastoral
reasons and for the sake of concelebration; and that \latin{Mysterium fidei}, drawn out of the
context of Christ's words, opens the way to the people's acclamation.}

\begin{studyframe}
\latin{\ldots ut eidem Precationi tres novi Canones adderentur statuimus. Attamen sive ut
pastoralibus, quas nominant, rationibus consuleretur, sive ut concelebratio expeditius
procederet, iussimus verba dominica in qualibet Canonis formula una eademque esse. Itaque in
quavis Precatione eucharistica illa sic proferri volumus: supra panem: \emph{Accipite et
manducate ex hoc omnes: Hoc est enim Corpus meum, quod pro vobis tradetur}; et supra calicem:
\emph{Accipite et bibite ex eo omnes: Hic est enim calix Sanguinis mei novi et aeterni
testamenti, qui pro vobis et pro multis effundetur in remissionem peccatorum. Hoc facite in meam
commemorationem.} Verba autem \emph{Mysterium fidei}, de contextu verborum Christi Domini
deducta, atque a sacerdote prolata, ad fidelium acclamationem veluti aditum aperiunt.}
\end{studyframe}

The comparison with the older Roman form is exact and checkable. The 1962 \latin{Missale
Romanum} prints, at the consecration of the chalice: \latin{Hic est enim Calix Sanguinis mei,
novi et aeterni testamenti: mysterium fidei: qui pro vobis et pro multis effundetur in
remissionem peccatorum}, followed by \latin{Haec quotiescumque feceritis, in mei memoriam
facietis}.\endnote{\latin{Missale Romanum}, editio typica 1962, \latin{Canon Missae}, marginal
nn.\ 1103--1105, printed p.\ 307, read visually at a rendered page image of the exact registered
facsimile on 2026-07-25. The reading was checked against the page image rather than the OCR.}
Three changes follow from the constitution's decision. \latin{Mysterium fidei} is removed from
the consecratory form and becomes the priest's cue for the people's acclamation, placed after
the elevation. The concluding command becomes \latin{Hoc facite in meam commemorationem},
conforming the Latin to 1~Corinthians 11:24--25 and Luke 22:19. And \latin{quod pro vobis
tradetur} is added to the form over the bread, which the older Roman form lacked.

Each of these can be described neutrally. The removal of \latin{mysterium fidei} from the form
takes out of the Lord's own words a phrase that is not in any Gospel account and restores it as
what it functionally is --- a statement about the mystery, spoken by the priest, to which the
people answer. The changed command aligns the Latin with the Pauline and Lukan text. The added
clause over the bread is likewise a conformation to the scriptural text. What is not neutral,
and should not be pretended otherwise, is that these are changes to the words of consecration
in the Roman Rite, made by papal act, on stated pastoral and practical grounds, and that they
are the deepest textual changes in the whole Order of Mass. A reader who wants the strongest
version of the traditionalist objection will find it here rather than in the offertory: the
form of consecration in the Latin West had stood substantially unchanged for well over a
millennium, and it was changed. The reply that carries weight is the one the constitution
itself gives --- that the dominical words are the same in all four prayers, taken from
Scripture, and that the mystery-of-faith phrase was not lost but repositioned --- and a reader
must weigh the two.

\subsection{The rubrical common shape (nn.\ 147--150)}

The Instruction adds the ceremonial that the four prayers share. The priest begins with hands
extended; he holds his hands extended over the offerings at the epiclesis; he pronounces the
words of the Lord clearly and distinctly, as the nature of those words requires; he shows the
consecrated Host and the chalice to the people and genuflects in adoration; and at the doxology
he takes the paten with the Host and the chalice and raises both.\endnote{The ceremonial
directions are printed in the Order of Mass itself at each prayer, most compactly at nn.\ 88--91
for the Roman Canon and in the parallel places of the other three. The formula ``In formulis
quae sequuntur, verba Domini proferantur distincte et aperte, prouti natura eorundem verborum
requirit'' stands identically before each institution narrative.} One of these rubrics was
touched by the 2008 reprint: at \latin{Institutio generalis} n.\ 149 the sequence of names in a
bishop's commemoration was reordered, so that when a bishop celebrates outside his own diocese
he now adds his brother bishop of that place before naming himself, where the 2002 text had the
order reversed.\endnote{\latin{Variationes} in \emph{Notitiae} 44 (2008) 368--369, the entry for
p.\ 48 of the \latin{Institutio generalis}. The change is one of the three submitted for papal
approval under Decree N.\ 652/08/L of 8 June 2008.}

\subsection{Which prayer, and when (\latin{IGMR} 365)}

The Instruction gives norms for choosing among the four, and they are more directive than is
often supposed.

\begin{itemize}
\item \textbf{Prayer I, the Roman Canon}, may always be used. It is more suitably said on days
assigned a proper \latin{Communicantes}, or at Masses enriched with a proper \latin{Hanc
igitur}, and in celebrations of the Apostles and of the Saints mentioned in it; likewise on
Sundays, unless for pastoral reasons the third prayer is preferred.
\item \textbf{Prayer II}, on account of its particular features, is more suitably used on
weekdays or in particular circumstances. Though it has its own preface, it may be used with
others, especially those that present the mystery of salvation in summary form.
\item \textbf{Prayer III} may be said with any preface, and its use is to be preferred on
Sundays and feasts.
\item \textbf{Prayer IV} has an unchangeable preface and gives a fuller summary of salvation
history. It may be used when a Mass has no proper preface, and on Sundays in Ordinary
Time.\endnote{\latin{IGMR} 365 a--d, read in the Latin. Each of the first three also carries a
provision for Masses for the dead: Prayer II has a special formula printed before \latin{Memento
etiam}; Prayer III has one inserted after \latin{Omnes filios tuos ubique dispersos}; Prayer IV
admits none, by reason of its structure.}
\end{itemize}

Two consequences follow that bear on ordinary parish practice. Prayer II is the shortest and
is, by the Instruction's own norm, the weekday prayer; its use as the standard Sunday prayer is
contrary to a stated norm, not merely to taste. And Prayer IV cannot be used when the day has a
proper preface, which excludes it from most solemnities and from the strong seasons; its
apparent rarity is largely a consequence of the calendar rather than of neglect.
